% ============================================================
% COMPOSER'S NOTE § 5 — How to Hear This
% ============================================================
% Reading guide. Names the three movements and their tempos:
%   — Geological (Ground)
%   — Chemical (Hold)
%   — Human (Touch)
% Describes the arc: subliminal → emerging → explicit;
% felt → witnessed → participated.
% Notes that global theories appear in service of specific
% material throughout — summoned by the material, not
% front-loaded in a prior chapter.
% Notes that the Coda names what was discovered (MSC,
% transmissible knowledge) rather than what was predetermined.
% Invites the reader to listen, not just to read.
% ============================================================

\section*{How to Hear This}

Three Movements. Three instruments. Three tempos.

\textbf{Movement I: Ground} moves at a geological pace---the tempo of sedimentation, of burial, of what the earth holds without announcing it. \textit{Buffalo Soldier} is the instrument: a wooden box filled with soil and rock, embedded with piezoelectric sensors. The layers of this Movement proceed from the surface (the problem of non-neutral tools) through the ground as archive, through encounter without archive, through the right to opacity, through the building of the instrument, to what the ground finally hears. The stakes arrive slowly, the way pressure does.

\textbf{Movement II: Hold} moves at a chemical pace---the tempo of dissolution, of transformation in suspension, of what water carries across time and space. \textit{Did Sun Ra Fear Cryosleep?} is the instrument: a vessel of water, hydrophones, dissolving artifacts. The logs of this Movement follow the wake, what water holds, temporal escape, the building of the instrument, and the dissolution that follows. This is the middle passage held open as a question.

\textbf{Movement III: Touch} moves at a human pace---the tempo of the body, of gesture and contact, of the politics of who gets to put their hands where. \textit{Don't Play with My Hair} is the instrument: a synthesis system that reads hair structure as data, that makes sound from what has been reached for and refused. The sections of this Movement trace signifying practice, the technical pipeline, haptic epistemology, braiding, the building of the instrument, and the final interruption.

Across all three Movements, the literature arrives when the material demands it. You will not find a prior chapter that assembles the theoretical sources before the work begins. Instead, post-phenomenological accounts of mediation appear when the earth vibrates in response to touch. Afrofuturist frameworks emerge when water becomes archive. Scholarship on embodiment and Black hair enters when the synthesis system confronts its own failure to see. Theory here is not scaffolding erected before the building. It is material---summoned, used, and answered to by the practice that called it.

The Coda names what was found. Not what was planned. It speaks to the transmissible knowledge that emerged from the making---the patterns, the failures, the procedures that can travel---and opens toward what this work cannot yet close.

The invitation is to listen, not just to read. Something has been sounding all along.

\adnote{The ``MSC'' referenced in the scaffolding (transmissible knowledge / what the Coda names) needs clarification---what does that acronym refer to? Confirm before the check-in. Also: the arc ``subliminal $\to$ emerging $\to$ explicit'' and ``felt $\to$ witnessed $\to$ participated'' from the scaffold are strong and could be woven in more explicitly here.}
